\documentclass[11pt]{article}

\usepackage[a4paper,margin=27mm]{geometry}
\usepackage{amsmath,amssymb,amsthm,mathtools}
\usepackage{enumitem}
\usepackage{hyperref}
\usepackage{xcolor}

\hypersetup{
  colorlinks=true,
  linkcolor=blue!50!black,
  citecolor=blue!50!black,
  urlcolor=blue!50!black
}

\theoremstyle{definition}
\newtheorem{definition}{Definition}[section]
\newtheorem{axiom}[definition]{Axiom}
\newtheorem{protocol}[definition]{Protocol}
\newtheorem{example}[definition]{Example}
\newtheorem{remark}[definition]{Remark}

\theoremstyle{plain}
\newtheorem{proposition}[definition]{Proposition}

\newcommand{\Eth}{\mathsf{Ethon}}
\newcommand{\Obs}{\mathsf{Obs}}
\newcommand{\Place}{\mathsf{Place}}
\newcommand{\Erase}{\mathsf{Erase}}
\newcommand{\Hide}{\mathsf{Hide}}
\newcommand{\Res}{\mathsf{Res}}
\newcommand{\Lic}{\mathsf{Lic}}
\newcommand{\Pol}{\mathsf{Pol}}
\newcommand{\Type}{\mathsf{Type}}
\newcommand{\Stab}{\mathsf{Stab}}
\newcommand{\Turn}{\mathsf{Turn}}
\newcommand{\Comp}{\mathsf{Comp}}
\newcommand{\Ex}{\mathsf{Ex}}
\newcommand{\Up}{\mathsf{Up}}
\newcommand{\Down}{\mathsf{Down}}
\newcommand{\Sdir}{\mathsf S}
\newcommand{\Tdir}{\mathsf T}
\newcommand{\CE}{\mathsf{CE}}
\newcommand{\UD}{\mathsf{UD}}
\newcommand{\ST}{\mathsf{ST}}
\newcommand{\Arg}{\mathsf{Arg}}
\newcommand{\FreeEth}{\mathsf{FreeEthon}}
\newcommand{\Fluc}{\mathsf{Fluc}}
\newcommand{\Rot}{\mathsf{Rot}}
\newcommand{\Pole}{\mathsf{Pole}}
\newcommand{\Npole}{\mathsf N}
\newcommand{\Spole}{\mathsf S}
\newcommand{\Spin}{\mathsf{Spin}}
\newcommand{\Bond}{\mathsf{Bond}}
\newcommand{\Overlap}{\mathsf{Overlap}}
\newcommand{\Qflav}{\mathsf{QFlav}}
\newcommand{\Qint}{\mathsf{QInt}}
\newcommand{\Sig}{\mathsf{Sig}}
\newcommand{\USD}{\mathsf{USD}}
\newcommand{\Fuse}{\mathsf{Fuse}}
\newcommand{\Fiss}{\mathsf{Fiss}}
\newcommand{\Plan}{\mathsf{Plan}}
\newcommand{\Civ}{\mathsf{Civ}}
\newcommand{\Law}{\mathsf{Law}}
\newcommand{\LawRev}{\mathsf{LawRev}}
\newcommand{\Disc}{\mathsf{Disc}}
\newcommand{\Disreg}{\mathsf{Disreg}}
\newcommand{\Ptr}{\mathsf{Ptr}}
\newcommand{\Quanta}{\mathsf{Quanta}}
\newcommand{\Rel}{\mathsf{Rel}}
\newcommand{\Read}{\mathsf{Read}}
\newcommand{\GenMes}{\mathsf{GenMes}}
\newcommand{\Ihara}{\mathcal I}
\newcommand{\Dir}{\mathcal D}
\newcommand{\Cat}{\mathbf{Cat}}
\newcommand{\defeq}{\mathrel{:=}}
\DeclareMathOperator{\supp}{supp}
\DeclareMathOperator{\diag}{diag}
\DeclareMathOperator{\Hom}{Hom}
\DeclareMathOperator{\SO}{SO}
\DeclareMathOperator{\Skew}{skew}

\title{Ethon -- A Pointer-Field Quantum Candidate}
\author{Shuhei Ihara}
\date{\today}

\begin{document}
\maketitle

\begin{abstract}
This paper introduces the ethon as a hypothetical quantum of a
boundary-indexed pointer field.  The field is not proposed as a
replacement for quantum mechanics, quantum field theory, QCD, relativity,
or established meson physics.  It is a typed state-transition model whose
local excitation, called an ethon, points from an unread relation to a
stable readable record.  The construction is inspired by Yukawa's act of
introducing a particle-like mediator for an unread nuclear relation, but
the ethon is not identified with the pion and is not asserted here to be
an experimentally confirmed particle.  Its core data are six directional
arguments, three rotational axes, boundary placement, hidden polarity,
residue, and bonding by complete overlap.  At a meson boundary, existing
QCD quark--antiquark labels may be read as ternary pointer-field
signatures in \(\{-1,0,1\}^{6}\); this is a readability hypothesis, not a
derivation of quark flavor or a replacement for QCD quantum numbers.  The
Ihara equation is presented as a typed update rule for the pointer-field
state, not as a physical equation of motion.  Physical interpretation is
therefore treated as a later research program: an ethon can be promoted
from formal quantum candidate to physical particle only after a field
action, covariance properties, spin-statistics assignment, couplings,
conservation laws, observables, and experimental constraints are supplied.
\end{abstract}

\tableofcontents

% -----------------------------------------------------------------------------
\section{Pointer Field Postulate}
% -----------------------------------------------------------------------------

\begin{definition}[Boundary-indexed pointer field]
For a typed space \(X\) and an admissible observer boundary \(b\), a
boundary-indexed pointer field is a formal field
\[
  \Ptr_b(X)
\]
whose local readable excitation points from an unread relation to a
stable readable record:
\[
  \Rel_b(a)\longrightarrow\Read_b(a).
\]
The field is defined at the level of typed state transitions.  A physical
field theory for \(\Ptr_b(X)\) is a further structure, not assumed at the
outset.
\end{definition}

The ethon is specified here as a quantum candidate of this pointer field.
Its existence is a primitive postulate of the present formal theory, not
a claim that an already confirmed natural particle has been detected.
The question of how an ethon would be produced or detected by an external
physical apparatus is part of the later physical research program.

\begin{axiom}[Ethon quantum-candidate existence]
For every typed space \(X\) and every admissible observer boundary \(b\),
there is a class of ethon candidates
\[
  \Eth_b(X)\subseteq \Quanta(\Ptr_b(X)).
\]
These candidates are formal objects of the implementation.  They may
be visible, hidden, or residue-carried at a boundary, but their physical
realization requires an additional target-domain license.
\end{axiom}

The interface is intentionally conservative.  It does not override quantum
mechanics, quantum field theory, relativity, meson physics, or ordinary
symbolic computation.  It states the additional structure that an ethon
candidate carries while remaining backward-compatible with those target
theories.

\begin{remark}[Yukawa inspiration]
Yukawa's meson paper is used here as a source of formal inspiration:
an unread interaction may become tractable by introducing a
particle-like mediator \cite{Yukawa1935}.  The present paper does not
claim institutional supervision, personal endorsement, historical
intention, or a physical derivation from Yukawa's original theory.
\end{remark}

\begin{remark}[Name and pronunciation]
The candidate is named \emph{Ethon} and pronounced ``Ee-thon.''  The name
is chosen to sound like a simple particle name and to connect the initial
sound to the Ihara equation.  The name does not supply evidence for
physical existence.
\end{remark}

\begin{definition}[Meson-boundary flavor reading]
Let the six quark-flavor names be denoted by
\[
  \Qflav_6
  =
  \{c,d,b,s,t,u\}.
\]
At a licensed meson boundary, the ethon direction frame may carry a
flavor-reading dictionary
\[
  \phi_Q:\Arg_{\Eth}\to\Qflav_6
\]
given by
\[
  \phi_Q(C)=c,\qquad
  \phi_Q(D)=d,\qquad
  \phi_Q(E)=b,\qquad
  \phi_Q(S)=s,\qquad
  \phi_Q(T)=t,\qquad
  \phi_Q(U)=u.
\]
The corresponding pointer-field reading is
\[
  \Qint_{\Eth}(\Qflav_6)
  =
  (\Arg_{\Eth},\ V_{\CE}\oplus V_{\UD}\oplus V_{\ST},\ \Rot).
\]
This is not a claim that the ethon derives, replaces, or exhausts QCD
flavor.  It is a readability hypothesis: at a meson boundary, existing
quark--antiquark labels may be read as directional pointer-field roles
whose visible organization is constrained by three-axis rotational data.
\end{definition}

\begin{definition}[Meson-boundary spin reading]
In the ethon model, a meson boundary is not merely a container for two
already-readable labels.  It is a boundary at which a coupled spin
pattern may make two quark roles readable:
\[
  \Spin_b^{\mathrm{mes}}:
  \Arg_{\Eth}\times\Arg_{\Eth}\times\Rot_b
  \longrightarrow
  \Read_b(q_i\bar q_j).
\]
Thus the phrase ``the meson defines quarks by spin'' means the following
restricted statement: within a licensed ethon reading, the meson boundary
defines the readable quark and antiquark roles through its spin-coupled
pointer pattern.  It does not mean that spin and flavor are the same
physical degree of freedom, nor that QCD quarks are replaced.
\end{definition}

\begin{definition}[Ternary meson-boundary signature]
A meson-boundary reading state has an external ethon signature
\[
  \Sig_{\Eth}\in\{-1,0,1\}^{\Arg_{\Eth}},
\]
with components \(\sigma_i\) for \(i\in\{C,D,E,S,T,U\}\).  The values
are read as
\[
  \begin{aligned}
  +1&:\text{quark-role-facing direction},\\
  -1&:\text{antiquark-role-facing direction},\\
   0&:\text{neutral, unpolarized, or closed direction}.
  \end{aligned}
\]
Thus the signature layer contains
\[
  |\{-1,0,1\}^{6}|=3^6=729
\]
possible six-direction interaction signs.  For an open-flavor basis
state \(q_i\bar q_j\) with \(i\neq j\), choose \(d_i,d_j\in\Arg_{\Eth}\)
with \(\phi_Q(d_i)=q_i\) and \(\phi_Q(d_j)=q_j\), and set
\[
  \sigma_{d_i}=+1,\qquad
  \sigma_{d_j}=-1,\qquad
  \sigma_k=0\quad(k\neq d_i,d_j).
\]
For a hidden-flavor state \(q_i\bar q_i\), the external signature is
neutral in that direction, \(\sigma_{d_i}=0\), but the channel is not
absent.  It is a closed zero mode carried by the rotational, polarity,
and bonding parameters of the ethon.  Thus the ternary signature is the
flavor-reading skeleton, not the whole meson identity; color singlet
structure, spin, orbital angular momentum, total angular momentum,
parity, charge conjugation, mass, and decay channels remain target-domain
data.
\end{definition}

\begin{definition}[Pointer-field quantum candidate]
A pointer-field quantum candidate is a formal local excitation whose
primary observable form is not a mass point or a medium, but an oriented
transition:
\[
  \epsilon^{\Obs}_b:
  \text{unread relation}
  \longrightarrow
  \text{stable readable entity}.
\]
The arrow is part of the candidate's observable signature.  An ethon is
the pointer-field quantum candidate specified in this document.
\end{definition}

\begin{remark}[Not a medium]
The ethon is not a revival of mechanical ether.  It may be placed in a
space, but it is not a substance filling space.  It is a direction-bearing
candidate quantum of relation, transition, and stabilization.
\end{remark}

\begin{definition}[Pure-meson reading]
A pure-meson reading is a meson-inspired operation stripped of fixed
empirical boundary data and retaining only its pointer role:
\[
  \text{unread relation}
  \longrightarrow
  \text{particle-like mediator}
  \longrightarrow
  \text{stable readable interaction}.
\]
The ethon is not thereby identified with an empirical meson.  It is a
pointer-field quantum candidate whose meson-like reading is prior to any
particular mass, charge, decay channel, detector boundary, or historical
realization.
\end{definition}

% -----------------------------------------------------------------------------
\section{Primitive Interface}
% -----------------------------------------------------------------------------

\begin{definition}[Typed space]
A typed space is any domain in which objects, observations, or states can
be distinguished by type.  Examples include physical space-time, Hilbert
space, phase space, a configuration space, a graph, a category of records,
a symbolic context, a neural embedding space, or a social decision space.
The definition does not assume a global time coordinate.
\end{definition}

\begin{definition}[Ethon instance]
An ethon instance in a typed space \(X\) is a tuple
\[
  \epsilon
  =
  (a,\nu,R,p,\varpi,\rho,\kappa,\lambda).
\]
Here:
\begin{enumerate}[label=(\roman*),leftmargin=*]
  \item \(a\in X\) is its placement address or contextual anchor;
  \item \(\nu=(\nu_{\CE},\nu_{\UD},\nu_{\ST})\) is its
  three-axis directional frame;
  \item \(R\in\SO(3)\) is its rotational attitude in that frame;
  \item \(p\) is its pole state, either absent or oriented as an
  \(S\)-face and \(N\)-pole;
  \item \(\varpi\in\mathbb R^3\) is its angular velocity vector;
  \item \(\rho\) is its residue ledger;
  \item \(\kappa\) is its stability and compatibility contract;
  \item \(\lambda\) is its license for interpretation inside a target
  theory, such as physical measurement, symbolic reading, or computation.
\end{enumerate}
The ethon is already real in the implementation.  The license
\(\lambda\) does not create the ethon; it only specifies how the real
ethon may be read in a particular target theory.
\end{definition}

\begin{definition}[Three rotational axes]
The ethon has three basic rotational axes:
\[
  \nu_{\CE},\qquad \nu_{\UD},\qquad \nu_{\ST}.
\]
The component \(\CE\) is the compression--exposition axis: it records
rotation between making many relations readable as one and exposing the
cuts that made that reading possible.  The component \(\UD\) is the
upward--downward axis, but a newly placed ethon has no actual up--down
polarity until the \(T\)-direction receives an entity.  The component
\(\ST\) is the stabilization--turnover axis: it records rotation between
stabilizing an observed pointer and turning that pointer backward into an
origin, cause, or reversed reading.
\end{definition}

\begin{axiom}[Three-directional existence]
Every ethon has all three rotational axes.  A model may suppress or hide
an axis at a boundary, but it may not define an ethon with only one or
two true axes:
\[
  \epsilon\in\Eth(X)
  \quad\Rightarrow\quad
  \nu(\epsilon)\in
  V_{\CE}\oplus V_{\UD}\oplus V_{\ST}.
\]
\end{axiom}

% -----------------------------------------------------------------------------
\section{Placement and Removal}
% -----------------------------------------------------------------------------

\begin{axiom}[Observer placement]
Given an observer boundary \(b\), a typed space \(X\), and an admissible
address \(a\in X\), the observer may place an ethon:
\[
  \Place_b(a):1\to\Eth(X).
\]
Placement does not mean physical manufacture from nothing.  It means that
the observer has introduced a pointer-particle coordinate at which a
relation can become readable.
\end{axiom}

\begin{axiom}[No removal]
There is no total removal operation
\[
  \Erase:\Eth(X)\to 1
\]
that annihilates an ethon and its residue.  Every attempted erasure
factors through hiding or residue transfer:
\[
  \Erase(\epsilon)
  =
  \Hide(\epsilon)
  \quad\text{or}\quad
  \Res(\epsilon).
\]
Thus an ethon can be made non-visible at a boundary, but the operation
that made it readable cannot be made to have never occurred.
\end{axiom}

\begin{proposition}[Residue conservation]
If an observer removes an ethon from the visible layer, then the total
record is conserved as hidden state plus residue:
\[
  Q(\epsilon)
  =
  Q(\Hide(\epsilon))+Q(\Res(\epsilon)),
\]
for every quantity valuation \(Q\) admitted by the model.
\end{proposition}

\begin{proof}
This is not a theorem of external physics.  It is the implementation
policy that replaces annihilation.  Since \(\Erase\) is not primitive,
every visible disappearance factors through an operation that preserves
hidden state, residue, or both.
\end{proof}

% -----------------------------------------------------------------------------
\section{Motion Law}
% -----------------------------------------------------------------------------

An ethon can transition in its three-axis state.  The transition does not
require a global time variable.  It is indexed by observation events,
boundary updates, or typed comparison steps.  The fundamental
fluctuation is not an extra dimension and not a string-like embedding.
It is rotation in the three-axis frame of the ethon.

\begin{definition}[Directional argument frame]
The six directional arguments of an ethon are
\[
  \Arg_{\Eth}
  =
  \{C,D,E,S,T,U\}.
\]
At a licensed meson boundary, they may also be read as the six roles of
the flavor-reading dictionary
\(\phi_Q:\Arg_{\Eth}\to\Qflav_6\).
Each directional argument carries a ternary interaction sign
\[
  \sigma_i\in\{-1,0,1\}.
\]
They are read in three paired axes:
\[
  (C,E),\qquad (U,D),\qquad (S,T).
\]
A placement state is a vector
\[
  x=(x_C,x_D,x_E,x_S,x_T,x_U)\in\mathbb R_{\geq0}^{6}.
\]
The component \(x_i\) records that an entity has been placed in the
directional argument \(i\).  The rotational frame is the three-dimensional
space
\[
  V_b\cong\mathbb R^3
  =
  \langle u_{\CE},u_{\UD},u_{\ST}\rangle,
\]
and the rotational attitude of an ethon is \(R\in\SO(3)\).
\end{definition}

\begin{axiom}[Free-functor inertness]
An ethon placed only as a free functor has no filled directional
argument:
\[
  \FreeEth(a)
  \quad\Rightarrow\quad
  x=0.
\]
In this state it has no up--down polarity and does not fluctuate:
\[
  p=\bot,
  \qquad
  \varpi=0,
  \qquad
  \Fluc_b=I_3.
\]
The identity attitude \(R=I_3\) may be used for bookkeeping, but it is
not an observed up--down orientation.  Placement alone is not motion.
\end{axiom}

\begin{definition}[Entity argument placement]
For \(i\in\Arg_{\Eth}\), placing an entity \(q\) in the \(i\)-th
directional argument updates the placement state by
\[
  x_i\mapsto x_i+w_i(q),
\]
where \(w_i(q)>0\) is the argument weight supplied by the model.  The
same entity may be placed in two different directional arguments, because
the arguments type the role of the entity rather than its bare identity.
\end{definition}

\begin{axiom}[One-argument hidden motion]
If exactly one directional argument is filled, internal motion may begin
but is not observable as fluctuation.  In a linearized implementation,
\[
  h_b(x)=K_bx
\]
may be nonzero, but the visible rotational fluctuation is still the
identity:
\[
  |\supp(x)|=1
  \quad\Rightarrow\quad
  h_b(x)\neq 0
  \ \text{is allowed, but}\
  \Fluc_b=I_3.
\]
\end{axiom}

\begin{definition}[Three-axis rotational fluctuation]
The visible ethon fluctuation between observation steps \(n\) and \(n+1\)
is the relative rotation
\[
  \Fluc_b(\epsilon_n)
  \defeq
  R_{n+1}R_n^{-1}
  \in\SO(3).
\]
Let \(u_{TS}\in V_b\) be the unit axis selected by the oriented
\(T\to S\) structure, and let \([u]_\times\in\mathfrak{so}(3)\) denote
the skew matrix satisfying \([u]_\times v=u\times v\).  The rotational
generator is
\[
  \Omega_{TS,b}(x)
  =
  \mathbf 1_{x_Tx_S>0}\,\omega_b [u_{TS}]_\times,
  \qquad
  \omega_b>0.
\]
The rotation update is
\[
  R_{n+1}
  =
  \exp\!\bigl(\Delta_n\Omega_{TS,b}(x_n)\bigr)R_n.
\]
When \(x_Tx_S>0\), the angular speed is constant:
\[
  \|\varpi_{n+1}\|=\omega_b.
\]
\end{definition}

\begin{axiom}[T-induced polarity]
A newly placed ethon has no north--south pole state.  A \(T\)-direction
placement creates hidden polarity:
\[
  x_T=0
  \quad\Rightarrow\quad
  p=\bot,
\]
whereas
\[
  x_T>0
  \quad\Rightarrow\quad
  p=
  (\Spole\text{-face}\to T,\ \Npole\text{-pole}\to S).
\]
This is a compass-like state: the \(S\)-face turns toward the
\(T\)-direction, while the \(N\)-pole points toward \(S\).  If
\(x_S=0\), this polarity is hidden and the rotation is not directly
observable.
\end{axiom}

\begin{axiom}[S-triggered constant self-rotation]
When a \(T\)-polarized ethon receives an \(S\)-direction placement, the
visible fluctuation channel opens as constant-speed self-rotation:
\[
  x_T>0,\quad x_S>0
  \quad\Rightarrow\quad
  \varpi_{n+1}=\omega_bu_{TS},
  \qquad
  \|\varpi_{n+1}\|=\omega_b.
\]
The ethon rolls in its own three-axis frame.  The speed is constant at
the boundary \(b\), although the axis coupling may still perturb all
three rotational axes.
\end{axiom}

\begin{example}[Rotational activation experiment]
The minimal activation table is:
\[
\begin{array}{c|c|c|c}
\text{placement} & p & \varpi & \Fluc_b\\
\hline
x=0 & \bot & 0 & I_3\\
x_T>0,\ x_S=0
  & (\Spole\to T,\Npole\to S) & 0 & I_3\\
x_T>0,\ x_S>0
  & (\Spole\to T,\Npole\to S) & \omega_bu_{TS}
  & \exp(\Delta_n\omega_b[u_{TS}]_\times).
\end{array}
\]
Thus the ethon is not made visible by an extra dimension.  It becomes
visible by constant-speed rotation after the \(T\to S\) axis is
completed.
\end{example}

% -----------------------------------------------------------------------------
\section{The Ihara Equation}
% -----------------------------------------------------------------------------

The Ihara equation is the central typed update rule of the implementation.
It turns directional argument placement into hidden motion, hidden
polarity, constant-speed rotation, residue, and meson-boundary pointer
reading.  It is not asserted here to be a physical equation of motion.

\begin{definition}[Meson-boundary pointer reading]
A meson-boundary pointer reading at boundary \(b\) is a morphism in the
boundary
category \(\GenMes_b\):
\[
  \mu\in
  \Hom_{\GenMes_b}\bigl(\Rel_b(a),\Read_b(a)\bigr),
\]
where \(\Rel_b(a)\) is an unread relation anchored at \(a\), \(\mu\) is
the meson-like pointer update, and \(\Read_b(a)\) is the stable readable
interaction obtained at that boundary.  The ethon is the formal
pointer-field quantum candidate whose update realizes this
meson-boundary reading without first choosing a particular physical mass,
charge state, decay channel, or detector boundary.
\end{definition}

\begin{definition}[Directional category]
The directional category of the ethon is
\[
  \Dir_{\Eth}\in\Cat.
\]
Its objects are the six directional arguments:
\[
  \operatorname{Ob}(\Dir_{\Eth})
  =
  \{C,D,E,S,T,U\}.
\]
The distinguished generating morphism is the premise--stabilization arrow
\[
  \tau_{TS}:T\longrightarrow S.
\]
The identities and all admissible composites are inherited from the
typed model, but the Ihara equation is based on \(\tau_{TS}\).  Thus the
mathematical entry point of the theory is not an unordered pair
\(\{S,T\}\), but an oriented categorical arrow \(T\to S\).
\end{definition}

\begin{definition}[Ihara phase space]
At boundary \(b\), the local Ihara phase space is
\[
  \mathcal M_b
  =
  V_b\times P_b\times \SO(3)\times \mathcal O_b\times R_b,
  \qquad
  V_b=V_{\CE}\oplus V_{\UD}\oplus V_{\ST},
  \qquad
  P_b=\mathbb R_{\geq0}^{\Arg_{\Eth}}.
\]
The pole-state set is
\[
  \mathcal O_b
  =
  \{\bot,\;(\Spole\text{-face}\to T,\Npole\text{-pole}\to S)\}.
\]
An element is written \((\nu,x,R,p,\rho)\), where \(\nu\) is the
three-axis motion state, \(x\) is the six-argument placement state,
\(R\in\SO(3)\) is the rotational attitude, \(p\) is the pole state, and
\(\rho\) is the residue state.  The support number is
\[
  \sigma(x)\defeq|\supp(x)|.
\]
\end{definition}

\begin{definition}[\(T\to S\) compass operator]
Let \(e_i\) denote the standard basis vector of \(P_b\) associated with
direction \(i\).  The categorical arrow \(\tau_{TS}:T\to S\) has two
algebraic representatives.  On placement coordinates it has the
rank-one form
\[
  R_{TS}
  =
  e_Se_T^\mathsf{T}:P_b\to P_b,
  \qquad
  R_{TS}x=x_Te_S.
\]
On rotational coordinates it selects a unit axis
\[
  u_{TS}\in V_b,
  \qquad
  \|u_{TS}\|=1.
\]
The placement-sensitive \(T\to S\) pole transfer is
\[
  \Theta_{TS}(x)
  =
  x_SR_{TS}x
  =
  x_Tx_Se_S,
\]
while the pole state induced by \(T\) is
\[
  P_T(x)
  =
  \begin{cases}
    \bot, & x_T=0,\\
    (\Spole\text{-face}\to T,\Npole\text{-pole}\to S),& x_T>0.
  \end{cases}
\]
\end{definition}

\begin{definition}[\(T\to S\) rotational generator]
The visible fluctuation generator is the map
\[
  \Omega_b:P_b\to\mathfrak{so}(3),
  \qquad
  \Omega_b(x)
  =
  \chi_{TS}(x)\,\omega_b[u_{TS}]_\times,
\]
where
\[
  \chi_{TS}(x)=\mathbf 1_{x_T>0\ \mathrm{and}\ x_S>0},
  \qquad
  \omega_b>0.
\]
Here \([u]_\times v=u\times v\).  The visible fluctuation is the rotation
\[
  \Fluc_b(x)
  =
  \exp\!\bigl(\Delta_n\Omega_b(x)\bigr)\in\SO(3).
\]
If \(\chi_{TS}(x)=1\), then
\[
  \varpi_b(x)=\omega_bu_{TS},
  \qquad
  \|\varpi_b(x)\|=\omega_b,
\]
so the ethon performs constant-speed self-rotation.  If
\(\chi_{TS}(x)=0\), then \(\Omega_b(x)=0\) and the visible fluctuation is
the identity rotation.
\end{definition}

\begin{definition}[Activation gates]
Define the support, polarity, and rotation gates
\[
  \chi_1(x)=\mathbf 1_{\sigma(x)\geq1},
  \qquad
  \chi_T(x)=\mathbf 1_{x_T>0},
  \qquad
  \chi_{TS}(x)=\mathbf 1_{x_T>0\ \mathrm{and}\ x_S>0}.
\]
The first gate starts hidden motion.  The second creates hidden
compass-like polarity.  The third opens observable constant-speed
rotation.
\end{definition}

\begin{axiom}[Ihara equation]
Let
\[
  \epsilon_n=(a,\nu_n,R_n,p_n,\varpi_n,\rho_n,\kappa,\lambda)
\]
be an ethon instance at observation step \(n\), let
\[
  x_n\in P_b,
  \qquad
  \eta_n\in V_b
\]
be its directional placement and unresolved input at boundary \(b\).  The
Ihara operator is the map
\[
  \Ihara_b:
  \mathcal M_b\times V_b
  \longrightarrow
  \mathcal O_b\times\mathbb R^3\times\SO(3)\times V_b\times R_b
  \times
  \Hom_{\GenMes_b}\bigl(\Rel_b(a),\Read_b(a)\bigr)
\]
defined by
\[
  \boxed{
  \Ihara_b(\epsilon_n,x_n,\eta_n)
  =
  (p_{n+1},\varpi_{n+1},R_{n+1},\nu_{n+1},\rho_{n+1},\mu_{n+1})
  }
\]
with
\[
\begin{aligned}
  p_{n+1}
  &=
  P_T(x_n),
  \\
  \varpi_{n+1}
  &=
  \chi_{TS}(x_n)\,\omega_bu_{TS},
  \\
  R_{n+1}
  &=
  \exp\!\bigl(\Delta_n[\varpi_{n+1}]_\times\bigr)R_n,
  \\
  m_{n+1}
  &=
  A_b\nu_n+K_bx_n+L_b\varpi_{n+1}+\eta_n,
  \\
  \nu_{n+1}
  &=
  \chi_1(x_n)m_{n+1},
  \\
  \rho_{n+1}
  &=
  \rho_n+\Gamma_b(p_{n+1},R_n,R_{n+1},\nu_n,\nu_{n+1},x_n),
  \\
  \mu_{n+1}
  &=
  \Pi_b(a,p_{n+1},R_{n+1},\nu_{n+1},\rho_{n+1})
  \in
  \Hom_{\GenMes_b}\bigl(\Rel_b(a),\Read_b(a)\bigr).
\end{aligned}
\]
Here \(A_b\) is the non-factorizing three-axis transition matrix,
\(K_b:P_b\to V_b\) converts six directional argument placements into
three-axis hidden motion, \(P_T\) is the \(T\)-induced pole map,
\(u_{TS}\) is the categorical \(T\to S\) rotation axis,
\(L_b:\mathbb R^3\to V_b\) embeds angular velocity into motion,
\(\Omega_b\) is the
skew-symmetric rotation generator, \(\Gamma_b\) records residue with
\(\Gamma_b(\bot,I_3,I_3,0,0,0)=0\), and \(\Pi_b\) reads the resulting
state as a meson-boundary pointer:
\[
  \Rel_b(a)
  \xrightarrow{\ \mu_{n+1}\ }
  \Read_b(a).
\]
\end{axiom}

\begin{proposition}[Rotational threshold]
The Ihara equation gives
\[
  x_T=0
  \Rightarrow
  p_{n+1}=\bot,\quad
  \varpi_{n+1}=0,\quad
  R_{n+1}=R_n,
\]
\[
  x_T>0,\ x_S=0
  \Rightarrow
  p_{n+1}=(\Spole\to T,\Npole\to S),\quad
  \varpi_{n+1}=0,\quad
  R_{n+1}=R_n,
\]
and
\[
  x_T>0,\ x_S>0
  \Rightarrow
  p_{n+1}=(\Spole\to T,\Npole\to S),\quad
  \varpi_{n+1}=\omega_bu_{TS},\quad
  R_{n+1}=\exp(\Delta_n\omega_b[u_{TS}]_\times)R_n.
\]
Thus observable ethon fluctuation is exactly the constant-speed
three-axis rotation opened by the completed \(T\to S\) placement.
\end{proposition}

\begin{proof}
The three cases follow directly from \(P_T\), \(\chi_{TS}\), and the
rotation update in the Ihara equation.
\end{proof}

\begin{proposition}[Activation trichotomy]
The Ihara equation implies the three activation regimes required by the
ethon:
\[
\begin{array}{c|c|c|c}
\text{condition} & p_{n+1} & \varpi_{n+1} & R_{n+1}\\
\hline
x_T=0 & \bot & 0 & R_n\\
x_T>0,\ x_S=0 & (\Spole\to T,\Npole\to S) & 0 & R_n\\
x_T>0,\ x_S>0 & (\Spole\to T,\Npole\to S)
  & \omega_bu_{TS}
  & \exp(\Delta_n\omega_b[u_{TS}]_\times)R_n.
\end{array}
\]
Thus a free ethon is inert, a \(T\)-placed ethon is hiddenly polarized,
and a \(T,S\)-placed ethon is visibly rotating.
\end{proposition}

\begin{proof}
This is immediate from the polarity gate \(\chi_T\) and the rotation gate
\(\chi_{TS}\).  Other placements may start hidden motion through
\(\chi_1\), but they do not open the ethon fluctuation channel unless the
\(T\to S\) rotation condition is satisfied.
\end{proof}

\begin{proposition}[T--S consequence]
If the first placed entity is a \(T\)-direction premise and the second
placed entity is a unit \(S\)-direction stabilizer, then the Ihara
equation gives
\[
  p_{n+1}
  =
  (\Spole\to T,\Npole\to S),
\]
and
\[
  R_{n+1}
  =
  \exp(\Delta_n\omega_b[u_{TS}]_\times)R_n,
  \qquad
  \|\varpi_{n+1}\|=\omega_b.
\]
If the second placed entity is instead \(j\neq S\), the \(T\)-induced
polarity remains hidden and the constant-speed rotation is not opened.
\end{proposition}

\begin{definition}[Transition step]
Let \(n\) index observation steps rather than physical time.  A transition
step is written
\[
  \nu_{n+1}
  =
  T_b(\nu_n,\eta_n),
\]
where \(b\) is the current observation boundary and \(\eta_n\) is the
unresolved rotational input at that step.
\end{definition}

\begin{axiom}[No axis isolation]
Motion, polarity, or rotation along one ethon axis always perturbs the
others.  In linearized form,
\[
  \Delta\nu
  =
  K_b\nu+L_b\varpi+\eta,
  \qquad
  K_b=
  \begin{pmatrix}
    k_{\CE\CE} & k_{\CE\UD} & k_{\CE\ST}\\
    k_{\UD\CE} & k_{\UD\UD} & k_{\UD\ST}\\
    k_{\ST\CE} & k_{\ST\UD} & k_{\ST\ST}
  \end{pmatrix},
\]
where no row and no column is allowed to be purely diagonal under a
nondegenerate observation boundary.  Equivalently, the transition kernel
does not factor as
\[
  P(\nu'_{\CE},\nu'_{\UD},\nu'_{\ST},R'\mid\nu,R)
  \neq
  P_{\CE}P_{\UD}P_{\ST},
\]
unless the model explicitly marks the boundary as degenerate.
\end{axiom}

\begin{definition}[S-facing observation]
The directly observed ethon pointer faces the stabilizing direction:
\[
  \Obs_b(\epsilon):
  \text{unread relation}
  \longrightarrow_{\Sdir}
  \text{stable readable record}.
\]
At the pole layer this means
\[
  p=(\Spole\text{-face}\to T,\Npole\text{-pole}\to S).
\]
The \(\Tdir\)-direction is not erased.  It remains as the hidden premise
toward which the \(S\)-face turns, and as the risk of back-reading the
stable pointer into a hidden origin or illegitimate cause.
\end{definition}

\begin{axiom}[Timeless admissibility]
An ethon may exist in a typed space that has no intrinsic time flow.
Its transitions are ordered by comparison, observation, or rewriting,
not necessarily by physical time:
\[
  n\prec n+1
  \quad\text{means}\quad
  \text{next readable transition},
\]
not necessarily ``later in time.''
\end{axiom}

% -----------------------------------------------------------------------------
\section{Stability}
% -----------------------------------------------------------------------------

\begin{definition}[Stable ethon]
An ethon is stable at boundary \(b\) when repeated observation steps
select the same pointer type and the same rotation class up to tolerated
residue:
\[
  \Obs_b(\epsilon_n)
  \simeq
  \Obs_b(\epsilon_{n+1})
  \quad
  \text{and}
  \quad
  \varpi_{n+1}\simeq\varpi_n
  \quad
  \text{modulo }\rho_b.
\]
The stable object may be still or constantly rotating.  It is not
residue-free.  It is readable because the residue is bounded, ledgered,
and compatible with the boundary.
\end{definition}

\begin{definition}[Ethon bond]
An ethon bond, formerly called an \(S\)-molecular bond, is not a line
between two separated ethons.  It is complete overlap at a boundary:
\[
  \Bond_b(\epsilon_1,\epsilon_2)
  \quad\Longleftrightarrow\quad
  \Overlap_b(\epsilon_1,\epsilon_2).
\]
Complete overlap means equality of the readable ethon state at that
boundary:
\[
  a_1=a_2,\quad
  x_1=x_2,\quad
  R_1=R_2,\quad
  p_1=p_2,\quad
  \Obs_b(\epsilon_1)=\Obs_b(\epsilon_2),
\]
with compatible residue ledgers.  Thus a bond is a two-ethon state in
which the two ethons occupy the same pointer, rotation, polarity, and
readable boundary condition.
\end{definition}

\begin{definition}[Ethon polymer]
An ethon polymer is a typed network
\[
  \mathcal P_{\Eth}
  =
  (Q_i,\epsilon_i,\beta_{ij},\kappa_{ij}),
\]
where \(Q_i\) are carriers, \(\epsilon_i\) are ethons, \(\beta_{ij}\) are
ethon bonds satisfying \(\Bond_b(\epsilon_i,\epsilon_j)\) when active,
and \(\kappa_{ij}\) are compatibility contracts.  The polymer is stable
when the network remains readable under the three-axis rotational
transition law.
\end{definition}

\begin{remark}[USD axis for polymers]
Earlier \(\USD\) fusion and \(\USD\) fission records are read here as
boundary-visible phases of ethon polymers.  This is a conservative
typing bridge: a \(\USD\) fission record does not by itself assert a
physical nuclear fission, chemical rupture, semantic failure, legal
crime, or policy fact.  Such descent still requires a target-domain
license.
\end{remark}

\begin{definition}[Rotational differential]
Let
\[
  \mathcal P_{\Eth,n}
\]
be an observed ethon polymer path at boundary \(b\), with rotational
attitudes \(R_n\in\SO(3)\) and observation increments \(\Delta_n>0\).
The three-axis rotational differential is
\[
  \dot R_n
  =
  \frac{1}{\Delta_n}
  \log(R_{n+1}R_n^{-1})
  \in\mathfrak{so}(3),
\]
whenever a boundary-compatible logarithm branch is available.  A
rotational discontinuity is recorded by
\[
  \Disc_b(n)=1
\]
when either no such branch is licensed, or
\[
  \left\|\dot R_{n+1}-\dot R_n\right\|>\theta_b,
\]
where \(\theta_b\) is the discontinuity tolerance of the boundary.  The
value \(\Disc_b(n)=0\) means only that no discontinuity is visible at the
chosen boundary resolution.
\end{definition}

\begin{definition}[Law reversal]
A law in the present paper is any licensed transition rule
\[
  \Law_b:
  \mathcal P_{\Eth,n}
  \longrightarrow
  \mathcal P_{\Eth,n+1}
\]
inside a target boundary \(b\).  The target may be mathematical,
semantic, physical, legal, institutional, or symbolic, but the target
reading requires its own license.  Law reversal occurs when a transition
is read as a lawful continuous update while a visible discontinuity is
left outside the residue ledger:
\[
  \begin{aligned}
  \LawRev_b(n)
  \quad\Longleftrightarrow\quad
  &\Law_b(\mathcal P_{\Eth,n}\to\mathcal P_{\Eth,n+1})
    \ \text{is asserted},\\
  &\Disc_b(n)=1,\qquad
    \Disc_b(n)\not\preceq\rho_{n+1}.
  \end{aligned}
\]
Thus a law reversal is not a moral verdict.  It is a typing failure in
which the \(T\)-side premise of rupture is read backward as if it were an
\(S\)-side stabilization.
\end{definition}

\begin{axiom}[Observer fusion boundary]
At the ethon layer, the admissible observer operation on polymers is
fusion:
\[
  \Fuse_b:
  \mathcal P_{\Eth,b}\times\mathcal P_{\Eth,b}
  \longrightarrow
  \mathcal P_{\Eth,b}.
\]
Fusion means that the observer places or bonds ethons so that one stable
composite polymer becomes readable at boundary \(b\).  There is no
primitive observer operation
\[
  \Fiss_b:
  \mathcal P_{\Eth,b}
  \longrightarrow
  \mathcal P_{\Eth,b}\times\mathcal P_{\Eth,b}
\]
that is licensed as an ethon operation by itself.
\end{axiom}

\begin{definition}[Observed polymer fission]
Observed polymer fission is not a primitive ethon operation.  It is the
boundary-visible record that a previously readable polymer is now read
through a discontinuity:
\[
  \Fiss_b^{\Obs}(\mathcal P_{\Eth})
  =
  \Obs_b\bigl(
  \mathcal P_{\Eth}
  \rightsquigarrow
  \mathcal P_{\Eth}^{(1)}\sqcup\mathcal P_{\Eth}^{(2)}
  \sqcup\rho
  \bigr).
\]
A physical nuclear fission, a molecular rupture, a semantic break in a
paper, or a legal offense may be mapped into this record only after a
separate target-domain license.
\end{definition}

\begin{definition}[Fission disregard]
Fission disregard is the record that a discontinuity has become visible
but has not been entered into the residue ledger:
\[
  \Disreg_b^{\Fiss}(n)
  \quad\Longleftrightarrow\quad
  \Disc_b(n)=1
  \ \text{and}\
  \Disc_b(n)\not\preceq\rho_{n+1}.
\]
In this terminology, a fission act is treated neutrally as disregard of
the fission discontinuity.  Its ethical, legal, physical, or scholarly
status is supplied only by the relevant target theory.
\end{definition}

\begin{proposition}[Law-reversal lemma]
Let
\[
  F=\Fuse_b(\mathcal P_1,\mathcal P_2)
\]
be a proposed observer operation, and let \(R_n\) be the resulting
three-axis rotational path.  If the operation is read as a continuous
fusion while its rotational differential has an unledgered
discontinuity,
\[
  \Disreg_b^{\Fiss}(n),
\]
then \(F\) is not classified as pure fusion.  It is retyped as fission
disregard and as a law reversal:
\[
  \Plan_b\bigl(F\mid\Disreg_b^{\Fiss}(n)\bigr)
  \in
  \Disreg_b^{\Fiss},
  \qquad
  \LawRev_b(n).
\]
\end{proposition}

\begin{proof}
By the observer fusion boundary, an observer may introduce a stable
composite polymer only through fusion.  By observed polymer fission,
fragmentation is not an ethon operation but a boundary-visible
discontinuity record.  If the proposed fusion is read as continuous while
\(\Disc_b(n)=1\) and that discontinuity is not included in
\(\rho_{n+1}\), then the operation has treated a \(T\)-direction turnover
as if it were already an \(S\)-facing stabilization.  This is precisely
the law reversal condition.  Hence the operation is ledgered as fission
disregard, not as pure fusion.
\end{proof}

\begin{remark}[Neutral target readings]
At a semantic boundary, law reversal may appear as a paper whose claims
ignore a discontinuity in their own transition rule.  At a legal
boundary, it may appear only if the legal system licenses the transition
as an offense.  At a physical boundary, it may appear as rupture or
fission only if the physical model supplies that descent.  The common
ethon record is merely
\[
  \Disreg_b^{\Fiss}(n)
  \quad\text{and}\quad
  \LawRev_b(n).
\]
It does not identify scholarly error, physical rupture, and crime with
one another.
\end{remark}

\begin{remark}[Matter hypothesis]
The strongest physical reading is the matter hypothesis:
\[
  \text{matter}
  =
  \text{stable observable phase of admissible ethon polymers}.
\]
This is a target implementation condition for a future physical reading
of ethon candidates as a matter theory.  It does not by itself replace
the established empirical content of existing physics.
\end{remark}

% -----------------------------------------------------------------------------
\section{Applications}
% -----------------------------------------------------------------------------

\begin{example}[Quantum carrier pair]
Let \(Q_1,Q_2\) be quantum carriers.  An ethon bond between their pointer
states is not a force line across distance.  It is a complete overlap of
two ethons at a boundary:
\[
  (Q_1,Q_2,\Bond_b(\epsilon_1,\epsilon_2))
  \longmapsto
  \text{observed pair state}.
\]
If the pair is entangled, the \(\ST\)-component is especially dangerous:
stable correlation may be turned backward into an illicit signal or
hidden cause.  The ethon ledger forbids that turn unless a physical
license supplies it.
\end{example}

\begin{example}[Pion compatibility and meson-boundary reading]
The pion remains an ordinary physical meson.  Ethon theory does not
change pion mass, spin, charge states, decay channels, or QCD status.
The pion is a licensed physical object of meson theory.  The ethon
reading keeps only the pointer-field role that a meson boundary may play
before the reading is fixed to a particular physical boundary.  Thus
\[
  \pi\text{-meson}
  =
  \text{QCD meson}
  +
  \text{empirical boundary data},
\]
whereas
\[
  \Eth
  =
  \text{pointer-field quantum candidate with meson-boundary reading}.
\]
This is a compatibility claim, not a replacement claim.
\end{example}

\begin{example}[Meson family interaction reading]
Under the flavor-shadow map \(\phi_Q\), a quark--antiquark meson label
\[
  q_i\bar q_j
\]
is first assigned a ternary ethon signature
\(\Sig_{\Eth}(q_i\bar q_j)\) only after a licensed meson spin boundary
\(\Spin_b^{\mathrm{mes}}\) is chosen.  For \(i\neq j\), it is read as
\[
  \sigma_{d_i}=+1,\qquad
  \sigma_{d_j}=-1,\qquad
  \sigma_k=0\quad(k\neq d_i,d_j),
  \qquad
  \phi_Q(d_i)=q_i,\quad
  \phi_Q(d_j)=q_j.
\]
For \(i=j\), the external signature is a closed zero mode on \(d_i\); it
is separated from absence by the spin-boundary, rotational, polarity, and
bonding data.
The resulting family is then classified by the three-axis rotational
state, the \(T\)-induced polarity, the \(S\)-triggered self-rotation, and
any ethon bonding condition.  This does not replace the empirical meson
tables; it supplies a possible pointer-field reading of meson labels as
ternary six-direction signatures refined by a spin-coupled, three-axis
ethon frame.
\end{example}

\begin{example}[Symbol and token]
A visible symbol or language-model token can be treated as a carrier in a
typed symbolic space.  The ethon is the pointer state that makes the mark
readable as an entity under a boundary.  Two symbolic ethons are bonded
only when their readable pointer, rotation, polarity, and boundary state
completely overlap.  Removing the symbol from view does not erase its
residue from memory, context, or model state.
\end{example}

\begin{example}[Discontinuity and recovery]
Nuclear fission, molecular rupture, semantic failure, and legal offense
are not the same event.  Ethon theory treats them neutrally as possible
target readings of one abstract pattern only when a licensed boundary
finds an unledgered discontinuity in a polymer transition.  Recovery
means preserving the usable knowledge while changing the polymer so that
the discontinuity is either avoided, bounded, or explicitly entered into
residue.
\end{example}

% -----------------------------------------------------------------------------
\section{Backward Compatibility}
% -----------------------------------------------------------------------------

\begin{protocol}[Classical physics]
Classical mechanics, field theory, and thermodynamics are preserved as
large-scale boundary readings.  The ethon may refine the bookkeeping of
relation, pointer, and residue, but it may not invalidate successful
classical limits.
\end{protocol}

\begin{protocol}[Relativity]
The ethon is not allowed to introduce a preferred physical rest frame.
If an ethon model is used in relativistic physics, its placement and
transition laws must be compatible with the relevant covariance
requirements of the target theory.
\end{protocol}

\begin{protocol}[Quantum theory]
The ethon may be used as a pointer-field quantum candidate for quantum
events only if it preserves no-signaling, the Born rule, uncertainty
relations, and the experimentally supported predictions of quantum
mechanics or quantum field theory.
\end{protocol}

\begin{protocol}[Physical particle program]
Promoting an ethon from formal quantum candidate to physical particle
requires additional data not supplied by the core implementation:
\[
  \begin{gathered}
  \text{field action},\quad
  \text{covariance law},\quad
  \text{spin-statistics assignment},\\
  \text{couplings},\quad
  \text{conservation laws},\quad
  \text{observables}.
  \end{gathered}
\]
It must also specify how the pointer field interacts, if at all, with
standard-model fields or gravity, and it must produce constraints or
predictions that can be compared with experiment.  Without this program,
the ethon remains a quantum candidate of the formal pointer field.
\end{protocol}

\begin{protocol}[String theory and extra dimensions]
The ethon does not require a ten-dimensional ambient space, compactified
extra dimensions, or a superstring target theory in order to fluctuate.
Its core fluctuation is rotation in the three-axis ethon frame.  String
or superstring theory may still be used as a licensed target theory, but
it is not a prerequisite for the ethon implementation.
\end{protocol}

\begin{protocol}[Meson theory]
The ethon may generalize the meson-like act of introducing a pointer
between unread carriers.  It may not rewrite the established physical
content of pion and meson physics without an independent physical
license.
\end{protocol}

\begin{protocol}[Smart functor records]
Earlier smart functor records are preserved as observed ethon records.
The functor is the observed arrow; the ethon is the pointer-field quantum
candidate whose observation appears as that arrow.
\end{protocol}

% -----------------------------------------------------------------------------
\section{Failure Modes}
% -----------------------------------------------------------------------------

\begin{enumerate}[leftmargin=*]
  \item \textbf{Medium collapse.}  Treating the ethon as a spatial ether
  rather than a pointer-field quantum candidate.
  \item \textbf{Axis isolation.}  Moving only one directional component
  while pretending the other two components are unaffected.
  \item \textbf{Dimensional inflation.}  Explaining ethon fluctuation by
  adding unnecessary ambient dimensions instead of using three-axis
  rotation.
  \item \textbf{Removal fiction.}  Claiming that a placed ethon can be
  erased without hidden state or residue.
  \item \textbf{Bond distance error.}  Treating an ethon bond as a line
  between separated objects rather than complete overlap of two ethons.
  \item \textbf{Discontinuity disregard.}  Treating a polymer transition
  with a visible rotational discontinuity as if it were a continuous
  lawful update.
  \item \textbf{T-direction overread.}  Turning a stable pointer backward
  into an unlicensed origin, cause, intention, or substance.
  \item \textbf{Target overdescent.}  Treating a formal ethon candidate as
  a licensed object of a specific target theory without the required
  target-domain contract.
\end{enumerate}

% -----------------------------------------------------------------------------
\section{Core Implementation Summary}
% -----------------------------------------------------------------------------

The minimal implementation is the following interface:
\[
  \epsilon
  =
  (a,\nu_{\CE},\nu_{\UD},\nu_{\ST},R,p,\varpi,\rho,\kappa,\lambda),
\]
with:
\begin{enumerate}[leftmargin=*]
  \item ternary six-direction signatures
  \(\Sig_{\Eth}\in\{-1,0,1\}^{6}\) for meson-boundary flavor-reading
  signs;
  \item pointer-field quantum-candidate formulation;
  \item meson-boundary quark-role readings through spin-coupled
  signatures;
  \item observer placement in arbitrary typed spaces;
  \item no total erasure operation;
  \item a three-axis rotational frame;
  \item no initial up--down polarity;
  \item hidden \(T\)-induced compass polarity;
  \item constant-speed self-rotation when \(T\) and \(S\) placements are
  both active;
  \item the \(T\to S\)-based Ihara equation as a typed
  meson-boundary pointer update rule;
  \item non-factorizing inter-axis dynamics;
  \item S-facing observed pointer form;
  \item ethon bonding as complete overlap of two ethons;
  \item observer fusion as the only primitive polymer operation;
  \item rotational discontinuity detection through
  \(\dot R_n=\Delta_n^{-1}\log(R_{n+1}R_n^{-1})\);
  \item the law-reversal lemma: unledgered rotational discontinuity is
  fission disregard, not pure fusion;
  \item residue conservation;
  \item backward compatibility with established target theories;
  \item physical-particle interpretation only by explicit license and
  additional field-theoretic data.
\end{enumerate}

Thus the ethon is implemented as a pointer-field quantum candidate:
freely placeable in the formal model, non-removable, stable across typed
spaces, independent of global time flow, rotationally fluctuating in
three axes, and dynamically coupled across all three rotational axes.

\paragraph{Acknowledgement and AI-use disclosure.}
The author expresses gratitude to Professor Hideki Yukawa for the meson
idea that made this paper's pointer-field formulation conceivable.  The
author used OpenAI GPT-5.5 for drafting support, terminology comparison,
and organization of the present formulation.  The author remains
responsible for the final claims, wording, and errors.

\begin{thebibliography}{9}

\bibitem{Yukawa1935}
Hideki Yukawa.
\newblock On the Interaction of Elementary Particles. I.
\newblock \emph{Proceedings of the Physico-Mathematical Society of Japan},
17:48--57, 1935.
\newblock Reprinted in \emph{Progress of Theoretical Physics Supplement},
1:1--10, 1955.
\newblock \url{https://doi.org/10.11429/ppmsj1919.17.0_48}.

\end{thebibliography}

\end{document}
